\par\noindent\begin{minipage}{\linewidth}
% !TeX root = ../教材题库_全部条目_转写工作稿.tex
% Source: XJ-B1-C04-HW-P104-Y1259
% 第4章 / 4.1 实数指数幂和幂函数 / 教材 p.104 / PDF 第111页
% 已对照原页转写题干及公式；待教师复核；未制作答案
% 原题及相邻原解截图（核对用）：../题目原图/XJ-B1-C04-HW-P104-Y1259-01.webp
% 本文件只在转写完成后作为题干使用；OCR 条目尚不是可靠题干。
\begin{bankitem}{习题4.1 第6题}{第4章 / 4.1 实数指数幂和幂函数 / 教材 p.104 / PDF 第111页}{XJ-B1-C04-HW-P104-Y1259}{已对照原页转写题干及公式；待教师复核；未制作答案}
化简：
\subq{1}$\left(\pi^{\sqrt2+1}\right)^{\sqrt2-1}$；
\subq{2}$\left(\left(\sqrt3\right)^{\sqrt[4]{64}}\right)^{\sqrt{25}}$。
\end{bankitem}
\end{minipage}\par\medskip
